\setcounter{section}{2}
\section{Demanda, horizonte y suficiencia de generación}
Se proyectan viviendas base de 2017 al año de estudio 2026 usando 1\% anual. Usuarios potenciales = ceil(viviendas\ensuremath{\times}1.01\textasciicircum{}(2026-2017)). Las categoríasI/I/II, simultaneidad 0.80,60 luminarias y 14 kW de servicios son hipótesis sin padrón. El resultado no acredita nuevas conexiones ni su distribución secundaria.

\begingroup\small\setlength{\LTpre}{4pt}\setlength{\LTpost}{6pt}
\begin{longtable}{*{4}{L{\dimexpr\linewidth/4-2\tabcolsep\relax}}}
\toprule
\textbf{Localidad} & \textbf{Viviendas 2017} & \textbf{Usuarios potenciales} & \textbf{kW} \\ \midrule
\endfirsthead
\toprule
\textbf{Localidad} & \textbf{Viviendas 2017} & \textbf{Usuarios potenciales} & \textbf{kW} \\ \midrule
\endhead
Huancasayani Centro & 95 & 104 & 49.920 \\
Huancasayani Bellavista & 140 & 154 & 73.920 \\
Huancasayani Lince & 86 & 95 & 30.400 \\
\bottomrule\end{longtable}\endgroup

P\_\allowbreak{}doméstica =154.240 kW; P\_\allowbreak{}AP=60\ensuremath{\times}0.0541=3.246 kW; P\_\allowbreak{}servicios=14 kW. P\_\allowbreak{}total=171.486 kW. Con FP0.90: Q=P\ensuremath{\times}tan(acos0.9)=83.054kvar; S=P/0.9=190.540 kVA. La carga se concentra en T2 para calcular el transporte, sin inventar ramales de usuario.

\begingroup\small\setlength{\LTpre}{4pt}\setlength{\LTpost}{6pt}
\begin{longtable}{*{4}{L{\dimexpr\linewidth/4-2\tabcolsep\relax}}}
\toprule
\textbf{Oferta} & \textbf{Reserva Limbani} & \textbf{Necesidad10 kV} & \textbf{Margen} \\ \midrule
\endfirsthead
\toprule
\textbf{Oferta} & \textbf{Reserva Limbani} & \textbf{Necesidad10 kV} & \textbf{Margen} \\ \midrule
\endhead
150 kW & 0 kW & 183.33 & -33.33 \\
150 kW & 25 kW & 208.33 & -58.33 \\
150 kW & 50 kW & 233.33 & -83.33 \\
200 kW & 0 kW & 183.33 & 16.67 \\
200 kW & 25 kW & 208.33 & -8.33 \\
200 kW & 50 kW & 233.33 & -33.33 \\
\bottomrule\end{longtable}\endgroup

La oferta debe contrastarse con la curva de disponibilidad y con la carga horaria simultánea de ambas localidades. Elevar tensión reduce pérdidas; no compensa un déficit de potencia activa de la central. La proyección a 10/20 años se exporta con cargas, potencia de generación y kVA para verificar también la capacidad aparente a la altitud del emplazamiento.

